\documentclass[11pt]{article}

\usepackage[margin=1in]{geometry}
\usepackage{amsmath,amssymb,amsthm}
\usepackage{mathtools}
\usepackage{microtype}
\usepackage{enumitem}

\newtheorem{assumption}{Assumption}

\title{Stationary Treatment Selection for Difference-in-Differences}
\author{}
\date{}

\begin{document}

\maketitle

\section{Statistical Setup}

Let $G\in\{0,1\}$ denote treatment-group membership, where $G=1$ denotes units that receive treatment after time $t_0$. Let $Y_t(0)$ and $Y_t(1)$ denote untreated and treated potential outcomes at time $t$. We are interested in the average treatment effect on the treated,
\[
\tau(t)
:=
\mathbb E[Y_t(1)-Y_t(0)\mid G=1].
\]

For $t<t_0$, treatment has not yet occurred, so
\[
Y_t=Y_t(0)
\]
for both treated and untreated groups. For $t\ge t_0$, we observe $Y_t(1)$ for the treated group and $Y_t(0)$ for the untreated group, while the counterfactual distribution
\[
Y_t(0)\mid G=1
\]
is unobserved.

Standard difference-in-differences identifies this counterfactual by imposing parallel trends on the conditional mean of $Y_t(0)$. We instead impose stability of the mechanism selecting units into treatment as a function of their untreated potential outcomes.

Define
\[
\gamma(y;t)
=
\mathbb P[G=1\mid Y_t(0)=y].
\]

We define the relative odds of treatment selection between two individuals with untreated potential outcomes $y$ and $y'$ as
\[
w(y,y';t)
=
\frac{
\gamma(y;t)/(1-\gamma(y;t))
}{
\gamma(y';t)/(1-\gamma(y';t))
}.
\]

\begin{assumption}[Stationary Treatment Selection]
The function $w(y,y';t)$ is constant in $t$.
\end{assumption}

The assumption allows the overall probability of treatment to vary over time, but requires the relative relationship between untreated potential outcomes and treatment selection to remain stable.

Under Assumption 1, there exist a time-specific function $\alpha(t)$ and a time-invariant function $\theta(y)$ such that
\[
\operatorname{logit}\mathbb P[G=1\mid Y_t(0)=y]
=
\alpha(t)+\theta(y).
\]

\section{Identification}

Let
\[
P_t^1
=
\mathcal L(Y_t(0)\mid G=1)
\]
and
\[
P_t^0
=
\mathcal L(Y_t(0)\mid G=0).
\]

Bayes' rule implies
\[
\frac{dP_t^1(y)}{dP_t^0(y)}
=
\frac{\mathbb P[G=0]}{\mathbb P[G=1]}
\frac{
\mathbb P[G=1\mid Y_t(0)=y]
}{
\mathbb P[G=0\mid Y_t(0)=y]
}.
\]

Under stationary treatment selection,
\[
\frac{
\mathbb P[G=1\mid Y_t(0)=y]
}{
\mathbb P[G=0\mid Y_t(0)=y]
}
=
\exp(\alpha(t)+\theta(y)).
\]
Therefore,
\[
\frac{dP_t^1(y)}{dP_t^0(y)}
=
C_t\exp(\theta(y)),
\]
where $C_t$ is a constant that depends only on $t$.

Consequently, the untreated mean of the treated population can be written as
\[
\mathbb E[Y_t(0)\mid G=1]
=
\frac{
\mathbb E[\exp(\theta(Y_t))Y_t\mid G=0]
}{
\mathbb E[\exp(\theta(Y_t))\mid G=0]
}.
\]

For pre-treatment periods, $Y_t=Y_t(0)$ for both groups. Thus, repeated pre-treatment observations identify the time-invariant weighting function $\theta(\cdot)$. Under Assumption 1, the same weighting function can then be applied to the post-treatment control distribution to recover the counterfactual untreated mean for the treated population.

Hence, for $t\ge t_0$,
\[
\tau(t)
=
\mathbb E[Y_t\mid G=1]
-
\frac{
\mathbb E[\exp(\theta(Y_t))Y_t\mid G=0]
}{
\mathbb E[\exp(\theta(Y_t))\mid G=0]
}.
\]

\section{Relation to Parallel Trends}

Stationary treatment selection provides an alternative to the usual parallel-trends restriction. Parallel trends assumes
\[
\mathbb E[Y_t(0)-Y_s(0)\mid G=1]
=
\mathbb E[Y_t(0)-Y_s(0)\mid G=0].
\]

In contrast, stationary treatment selection allows the untreated potential-outcome distributions of the two groups to evolve differently over time, provided that their density ratio retains the form
\[
\log
\frac{dP_t^1(y)}{dP_t^0(y)}
=
\alpha(t)+\theta(y).
\]

Thus, parallel trends transports changes in outcomes across groups, while stationary treatment selection transports a stable reweighting function across time.

The assumption may be natural in settings where treatment assignment follows a persistent selection rule. For example, treatment capacity or the overall treatment rate may change over time, while the relative tendency to treat individuals with different untreated potential outcomes remains stable.

Multiple pre-treatment periods also provide a diagnostic for the assumption. Stationary treatment selection implies that
\[
\log\frac{dP_t^1(y)}{dP_t^0(y)}
-
\log\frac{dP_t^1(y')}{dP_t^0(y')}
\]
should remain approximately constant across pre-treatment periods. This provides an analogue of a pre-trends diagnostic for the stability of treatment selection.

\end{document}
